% Preface: the README's introduction, Om kursen and Efter kursen; the prerequisites, software,
% hardware and references from info/README.md; how a lecture is laid out and the kinds of
% exercise from lectures/README.md; and how a lab is done, with its safety rules, from
% labs/README.md. The course's schedule, dates and class are left out on purpose: the book is
% meant to outlive any one run of the course.

\chapter{Förord}
\markboth{Förord}{Förord}

Den här boken tar dig från ettor och nollor till en mikrodator som styr något i verkligheten.
Först talsystem, logiska grindar och boolesk algebra, byggda både med strömkontakter på 24~V och
med IC-kretsar på kopplingsdäck. Sedan vippan, den byggsten som ger en krets minne, och därifrån
mikrodatorn: en ATmega328P som programmeras i assembler, först i simulatorn och till sist på ett
riktigt kort där den styr lysdioder och läser av en knapp.

Bokens text är på svenska. All kod och alla kommentarer i koden är på engelska, vilket är det du
kommer att möta i datablad, instruktionslistor och på nätet.

\section*{Bokens tre delar}
\paragraph{Del 1: Digitala kombinatoriska nät, kapitel 1--5} Binära, hexadecimala och decimala tal
och omvandling mellan dem. Logiska grindar, sanningstabeller och boolesk algebra, och hur ett
uttryck förenklas, algebraiskt och med Karnaughdiagram. Samma logik realiseras på två sätt: med
tryckknappar och en lampa på 24~V i \textbf{Labb~1}, och med IC-kretsar ur 74-serien på
kopplingsdäck i \textbf{Labb~2}.

\paragraph{Del 2: Från vippa till mikrodator, kapitel 6} Sekvensnät och vippor: hur en krets kan
minnas. Register, räknare och minne är vippor, och mikrodatorn är register och minne kring en
aritmetisk enhet, byggd av de grindar del~1 handlade om. Kapitlet slutar i mikrodatorns
blockschema.

\paragraph{Del 3: Mikrodatorn och assemblerprogrammering, kapitel 7--11} ATmega328P programmeras i
assembler i Microchip Studio: talformat, aritmetik och flaggor, logiska operationer, villkorliga
hopp, loopar och tidsfördröjningar, in- och utportar, subrutiner och stacken. \textbf{Labb~3} löper
genom alla fem kapitlen, och avslutas med ett styrobjekt, ett trafikljus, som körs på ett Arduino
Uno-kort.

\paragraph{Repetition, kapitel 12} Hela kursen i sammandrag, med de vanligaste felen och övningar
i provets form, inför det skriftliga provet.

\section*{Det här ska du kunna efteråt}
När du har arbetat dig igenom boken ska du kunna:
\begin{itemize}
  \item Förklara grindar och logiska funktioner, och beskriva logisk algebra och förenkling av
    booleska uttryck.
  \item Beskriva det binära, det hexadecimala och det decimala talsystemet, och omvandla tal
    mellan dem.
  \item Syntetisera ett digitalt kombinatoriskt nät, med IC-kretsar och med kontakter, och använda
    Karnaughdiagram för att analysera och minimera det.
  \item Analysera ett digitalt kombinatoriskt nät: från krets till uttryck till sanningstabell.
  \item Redogöra för mikrodatorns uppbyggnad på blockschemanivå.
  \item Använda assemblerspråk för att programmera en mikrodator, och programmera den att styra ett
    styrobjekt.
\end{itemize}

\section*{Vad du behöver kunna innan}
Inga förkunskaper i elektronik eller programmering förutsätts. Det räcker med matematik på
gymnasienivå: potenser, och de fyra räknesätten utan miniräknare.

Kursen boken är skriven för går snabbt. Tolv pass om tre timmar rymmer fyra ämnen som var och ett
skulle kunna fylla en egen kurs, så \textbf{varje kapitel förutsätter att det föregående är läst
och att dess övningar är gjorda.}

\section*{Hur boken är upplagd}
Boken är satt från en kurs med tolv pass om tre timmar, och varje kapitel är ett pass. Kapitlets
avsnitt är passets appendix, i samma ordning och med samma underavsnitt, så att ett kapitel här
och ett pass i kursmaterialet kan läsas sida vid sida:

\begin{booktable}{@{}p{5.2em}p{3.4em}L@{}}
  \thead{Kapitel} & \thead{Pass} & \thead{Ämne} \\
  \midrule
  1 & \file{L01} & Talsystem och binära koder \\
  2 & \file{L02} & Logiska grindar, boolesk algebra och kontaktnät \\
  3 & \file{L03} & Analys och felsökning av kontaktnät, med \textbf{Labb~1} \\
  4 & \file{L04} & Karnaughdiagram och IC-kretsar \\
  5 & \file{L05} & Felsökning av grindnät, med \textbf{Labb~2} \\
  6 & \file{L06} & Sekvensnät, vippor och mikrodatorns uppbyggnad \\
  7 & \file{L07} & Mikrodatorn och Microchip Studio; \textbf{Labb~3} börjar \\
  8 & \file{L08} & Aritmetik, logik och flaggor \\
  9 & \file{L09} & Programflöde, loopar och tidsfördröjningar \\
  10 & \file{L10} & Inporten, subrutiner och stacken \\
  11 & \file{L11} & Styrobjekt på riktig hårdvara; \textbf{Labb~3} slutförs \\
  12 & \file{L12} & Repetition inför det skriftliga provet \\
\end{booktable}

Varje kapitel börjar med vad det innehåller, arbetar sig igenom materialet, sammanfattar vad du
nu ska kunna förklara, och slutar med sina övningar. Bilagorna håller det som används vid sidan
av kapitlen: de tre laborationshandledningarna, en guide till Microchip Studio, referensbladet med
AVR-instruktionerna, en beskrivning av det skriftliga provet med ett övningsprov och dess
lösningsförslag, och sist lösningsförslagen till övningarna.

\section*{Övningarnas sorter}
Varje övning är märkt med sin sort:

\begin{booktable}{@{}p{9em}L@{}}
  \thead{Sort} & \thead{Vad den ber dig göra} \\
  \midrule
  \kinds{Förståelse} & Förklara något från kapitlet med egna ord. Inga verktyg. \\
  \kinds{Räkna för hand} & Räkna ut något på papper: en omvandling, en förenkling, en
    cykelräkning. \\
  \kinds{Konstruktion} & Ta fram en lösning: ett grindnät, ett kontaktnät, en flödesplan. \\
  \kinds{Program} & Skriv assembler, och kör det i simulatorn. \\
  \kinds{Kontroll} & En per kapitel. Se nedan. \\
  \kinds{Fördjupning} & En övning för den som vill gå längre. Den kan hoppas över. \\
\end{booktable}

\textbf{Kontrolluppgiften är kursens signaturövning.} Du räknar ut ett svar för hand, du
kontrollerar det med ett verktyg, i CircuitVerse, i simulatorn eller på labbstationen, och du
förklarar varför de två stämmer, eller varför de inte gör det. Att förutsäga ett beteende och
sedan kontrollera det är den färdighet resten av kursen vilar på.

\section*{Lösningsförslagen}
\textbf{Lösningsförslagen till övningarna finns i bilaga~\ref{app:answers}}, sist i boken, kapitel
för kapitel. Där står svaret på varje övning som räknas, konstrueras eller förklaras, med
uträkningen, och ofta med de vanligaste felaktiga svaren och varför de är fel.

Två slags uppgifter har inga lösningsförslag i boken. \textbf{Programuppgifterna}, där du skriver
assembler: ett program lär du dig skriva genom att skriva det, och simulatorn visar om det gör vad
det ska. Delar av en programuppgift som räknas eller förklaras, en förutsägelse av vad registren
innehåller eller en cykelräkning, har däremot sina svar i bilagan. \textbf{Laborationerna}
redovisas för läraren vid stationen, och deras svar är det du visar där.

Försök alltid själv först: ett svar du har läst är mycket mindre värt än ett du har räknat fram
och sedan kontrollerat.

\section*{Laborationerna}
Kursen har tre laborationer, som görs i par. Labb~1 och Labb~2 görs på varsitt pass; Labb~3 löper
över de fem mikrodatorkapitlen. Handledningarna finns i bilaga~\ref{app:lab1} till
\ref{app:lab3}.
\begin{itemize}
  \item \textbf{Förbered dig hemma.} Varje handledning börjar med förberedelseuppgifter:
    sanningstabeller, uttryck och kopplingsscheman som ska vara klara när passet börjar. Utan dem
    hinner man inte klart, och det är den enskilt vanligaste orsaken till att en laboration inte
    blir godkänd på passet.
  \item \textbf{Arbeta i par}, och byt roll: den ena kopplar eller skriver, den andra kontrollerar
    mot schemat. Båda ska kunna förklara lösningen.
  \item \textbf{Förutsäg, sedan kontrollera.} Skriv ned vad du förväntar dig innan du slår på
    strömmen eller kör programmet. En avvikelse mellan förutsägelse och resultat är information:
    den säger att schemat, kopplingen eller programmet har ett fel, och att du nu vet var du ska
    leta.
  \item \textbf{Redovisa för läraren} varje uppgift som är markerad \kindtag[fillaccent]{Redovisa}.
    Läraren prickar av den i redovisningslistan.
\end{itemize}

\begin{warning}
\textbf{Säkerhet.}
\begin{itemize}
  \item Koppla \textbf{alltid} med spänningen frånslagen, och slå på först när kopplingen är
    kontrollerad.
  \item 24~V DC är ingen farlig spänning att beröra, men en kortslutning kan ge stora strömmar,
    heta sladdar och en bränd säkring i aggregatet. En glödlampa blir dessutom varm.
  \item Koppla aldrig 24~V till kopplingsdäcket eller till Arduino-kortet. IC-kretsarna och
    mikrodatorn ska matas med 5~V, och 24~V förstör dem på ett ögonblick.
  \item IC-kretsar vänds lätt fel. En krets som blir varm är felvänd eller felkopplad: bryt
    matningen direkt.
\end{itemize}
\end{warning}

\section*{Programvara}
\begin{itemize}
  \item \textbf{CircuitVerse}, \url{https://circuitverse.org/simulator}: en gratis,
    webbläsarbaserad logiksimulator, där grindnät byggs och testas innan de kopplas upp. Inget att
    installera.
  \item \textbf{Microchip Studio 7}, \url{https://www.microchip.com/en-us/tools-resources/develop/microchip-studio}:
    en gratis utvecklingsmiljö för AVR, där programmen i Labb~3 skrivs, assembleras och simuleras.
    Den installeras före kapitel~\ref{c7:ch} enligt bilaga~\ref{app:studio}. Microchip Studio finns
    bara för Windows; den som har en Mac använder skolans datorer eller en virtuell
    Windowsmaskin.
  \item \textbf{Arduino IDE}, \url{https://www.arduino.cc/en/software}: används en enda gång, i
    kapitel~\ref{c11:ch}, för att ta fram kommandot som programmerar kortet (se
    bilaga~\ref{app:studio}).
\end{itemize}
Under kursen går det bra att kontrollera omvandlingar med kalkylatorn i Windows, i läget
\textbf{Programmerare}, men räkna alltid för hand först. På provet finns ingen miniräknare.

\section*{Utrustning}
\textbf{Labb~1} görs vid en station med ett 24~V DC-aggregat, två 4-poliga tryckknappar med både
slutande (NO) och brytande (NC) kontakter, en glödlampa för 24~V, kopplingssladdar och en
multimeter för felsökning.

\textbf{Labb~2} görs vid ett kopplingsdäck med kringutrustning: 5~V-matning, strömbrytare eller
knappar för ingångarna och lysdioder för utgångarna. Till det kommer IC-kretsarna 74HC08 (AND),
74HC32 (OR) och 74HC04 (inverterare) till Labb~2-1, 74HC00 (NAND) och 74HC02 (NOR) till Labb~2-2,
kopplingstråd och en multimeter.

\textbf{Labb~3} görs med ett Arduino Uno-kort (ATmega328P) med USB-kabel, ett kopplingsdäck, tre
lysdioder (röd, gul och grön) med förkopplingsmotstånd på 220--330~Ω, en tryckknapp och
kopplingstråd. Det mesta av Labb~3 görs i simulatorn i Microchip Studio; kortet behövs först i
kapitel~\ref{c11:ch}.

\section*{Referenser}
Två dokument från tillverkaren används som referens i mikrodatordelen:
\begin{itemize}
  \item \emph{ATmega328P datasheet}:
    \url{https://ww1.microchip.com/downloads/en/DeviceDoc/Atmel-7810-Automotive-Microcontrollers-ATmega328P_Datasheet.pdf}
  \item \emph{AVR Instruction Set Manual}:
    \url{https://ww1.microchip.com/downloads/en/devicedoc/atmel-0856-avr-instruction-set-manual.pdf}
\end{itemize}
Det du behöver av dem under kursen, instruktionerna, flaggorna och cykeltiderna, finns samlat på
referensbladet i bilaga~\ref{app:avr}.
